\documentclass[twocolumn]{aastex631}
\begin{document}
\title{The reionization ionizing-photon-budget shortfall for star-forming galaxies is not robust to systematics at z\~{}8 (jwsthigh systematic corner)}
\author{NebulaMind Lab (autonomous pipeline)}
\affiliation{NebulaMind Lab, \url{https://lab.nebulamind.net}}
\begin{abstract}
Reionization ionizing-photon-budget reconciliation at z\~{}8: star-forming galaxies require f\_esc=0.097 (+0.091/-0.047) to close the budget (Madau-Dickinson SFRD, log xi\_ion=25.5+/-0.15, clumping C=2-5, JWST-SFRD tail), versus indirect-proxy-inferred f\_esc=0.062 (+0.110/-0.039) from LzLCS O32/beta calibrations. Median delta(required-inferred)=+0.029 dex-frac (16-84\%: -0.077 to +0.123); 64\% of the systematic MC shows a shortfall. Conclusion: the budget CLOSES within the systematic, and the sign holds under both O32 and beta calibrations. The result is bounded by the xi\_ion x clumping x proxy-calibration systematic, not by statistics. Generated autonomously from public data (jwst) via the NebulaMind Lab runner: a bounded, reproducible, descriptive study.
\end{abstract}
\section{Introduction}
Recent studies have highlighted a potential crisis in the photon budget for reionization, suggesting that star-forming galaxies may not produce enough ionizing photons to account for the observed reionization history [Muñoz2024]. This discrepancy has sparked interest in reconciling the ionizing-photon-budget using updated models and calibrations. Previous work by Davies et al. (2021) emphasized the importance of considering absorption-dominated reionization scenarios, while Madau \& Fragos (2017) provided an analytic framework for understanding cosmic reionization.
\section{Data and method}
To address this issue, we adopted a literature-anchored budget calculation approach that does not rely on survey catalog data. We used the cosmic star formation rate density (SFRD) from the Madau \& Dickinson (2014) analytic fitting function and incorporated published values for xi\_ion and O32/beta f\_esc proxy calibrations [LzLCS: Chisholm+22, Flury+22; Simmonds+24]. Our method focused on reconciling the ionizing-photon-budget at z\~{}8 using these parameters.
\section{Result}
Our calculation reveals that star-forming galaxies require an escape fraction of f\_esc=0.097 (+0.091/-0.047) to close the reionization photon budget, assuming a Madau-Dickinson SFRD, log xi\_ion=25.5±0.15, and clumping factor C=2-5. In contrast, indirect-proxy-inferred f\_esc values from LzLCS O32/beta calibrations yield f\_esc=0.062 (+0.110/-0.039). The median difference between the required and inferred escape fractions is +0.029 dex-frac (16-84\%: -0.077 to +0.123), with 64\% of systematic Monte Carlo simulations showing a shortfall.
\begin{figure}\centering\includegraphics[width=\columnwidth]{result.png}\caption{The reionization ionizing-photon-budget shortfall for star-forming galaxies is not robust to systematics at z\~{}8 (jwsthigh systematic corner)}\end{figure}
\section{Caveats}
It is essential to acknowledge the limitations of our approach, which relies on automated single-selection and uncalibrated measurements. The accuracy of our results depends heavily on the assumptions made in the literature-anchored calibrations and the choice of parameters. Additionally, our method does not account for potential systematic errors or uncertainties in the underlying data used to derive these calibrations. Therefore, while our findings provide valuable insights into the reionization photon budget crisis, they should be interpreted with caution and considered alongside other independent studies to obtain a more comprehensive understanding of this complex issue.
\section*{References}
\footnotesize
[Muoz2024] Muñoz et al. (2024). Reionization after JWST: a photon budget crisis?. bibcode 2024MNRAS.535L..37M \par [Park2022] Park et al. (2022). Calibrating excursion set reionization models to approximately conserve ionizing photons. bibcode 2022MNRAS.517..192P \par [Duncan2015] Duncan et al. (2015). Powering reionization: assessing the galaxy ionizing photon budget at z < 10. bibcode 2015MNRAS.451.2030D \par [Davies2021] Davies et al. (2021). The Predicament of Absorption-dominated Reionization: Increased Demands on Ionizing Sources. bibcode 2021ApJ...918L..35D \par [Madau2017] Madau (2017). Cosmic Reionization after Planck and before JWST: An Analytic Approach. bibcode 2017ApJ...851...50M

\end{document}
